% ----- CHAUPAI 23 -----

\newpage

\subsection*{\deva{त्रयोविंशतितम चौपाई} \(Twenty-Third Chaupai\)}
\addcontentsline{toc}{subsection}{\deva{त्रयोविंशतितम चौपाई} \(Twenty-Third Chaupai\) — \deva{आपन तेज...}}

\begin{minipage}[t]{0.55\textwidth}
\vspace{0pt}

{\large \linespread{1.5}\selectfont
\deva{आपन तेज सम्हारो आपै।}\\
\deva{तीनों लोक हाँक तें काँपै॥}
\par}

\vspace{0.5em}

{\large \linespread{1.5}\selectfont
\telu{ఆపన తేజ సమ్హారో ఆపై।}\\
\telu{తీనోం లోక హాఁక తేఁ కాఁపై॥}
\par}

\vspace{0.5em}

{\large \linespread{1.3}\selectfont
\textit{āpana teja samhāro āpai |}\\
\textit{tīnõ loka hā̃ka tẽ kā̃pai ||}
\par}
\end{minipage}%
\hfill
\begin{minipage}[t]{0.42\textwidth}
\vspace{0pt}
\begin{tcolorbox}[anchorbox]
\textbf{Self-Regulation:} He possessed immense, earth-shaking energy and talent, but he kept it under perfect control (\textit{Samharo Aapai}). Raw intelligence, talent, and energy are completely useless without the self-discipline to regulate and direct them properly.
\vspace{0.5em}\newline He completely regulated his blinding, destructive energy, using it only precisely when the situation demanded it.\newline\null\hfill\href{https://www.valmiki.iitk.ac.in/sloka?field_kanda_tid=5&language=dv&field_sarga_value=42}{\textit{Sundara Kanda, 42:34}}
\end{tcolorbox}
\end{minipage}

\vspace{0.5em}
\subsubsection*{\textbf{Visual Rhythm Map:}}
\begin{table}[h!]
\centering
\renewcommand{\arraystretch}{1.2}
\setlength{\tabcolsep}{0pt}
\begin{tabularx}{\textwidth}{|*{16}{>{\centering\arraybackslash\hsize=1.0\hsize}X|}}
\hline
\multicolumn{4}{|>{\centering\arraybackslash\hsize=4.0\hsize}X|}{\textbf{\guru{आ}\laghu{प}\laghu{न}} \par (4)} &
\multicolumn{3}{>{\centering\arraybackslash\hsize=3.0\hsize}X|}{\textbf{\guru{ते}\laghu{ज}} \par (3)} &
\multicolumn{5}{>{\centering\arraybackslash\hsize=5.0\hsize}X|}{\textbf{\laghu{स}\guru{म्हा}\guru{रो}} \par (5)} &
\multicolumn{4}{>{\centering\arraybackslash\hsize=4.0\hsize}X|}{\textbf{\guru{आ}\guru{पै}} \par (4)} \\
\hline
\end{tabularx}
\vspace{-1.5em}
\end{table}

\begin{table}[h!]
\centering
\renewcommand{\arraystretch}{1.2}
\setlength{\tabcolsep}{0pt}
\begin{tabularx}{\textwidth}{|*{16}{>{\centering\arraybackslash\hsize=1.0\hsize}X|}}
\hline
\multicolumn{4}{|>{\centering\arraybackslash\hsize=4.0\hsize}X|}{\textbf{\guru{ती}\guru{नों}} \par (4)} &
\multicolumn{3}{>{\centering\arraybackslash\hsize=3.0\hsize}X|}{\textbf{\guru{लो}\laghu{क}} \par (3)} &
\multicolumn{3}{>{\centering\arraybackslash\hsize=3.0\hsize}X|}{\textbf{\guru{हाँ}\laghu{क}} \par (3)} &
\multicolumn{2}{>{\centering\arraybackslash\hsize=2.0\hsize}X|}{\textbf{\guru{तें}} \par (2)} &
\multicolumn{4}{>{\centering\arraybackslash\hsize=4.0\hsize}X|}{\textbf{\guru{काँ}\guru{पै}} \par (4)} \\
\hline
\end{tabularx}
\vspace{-1.5em}
\end{table}

\vspace{0.5em}
\textbf{ \deva{सरल-संस्कृतम्}:}\\
 \deva{भवतः प्रचण्डं तेजः भवान् स्वयम् एव सम्भालितुं शक्नोति।}\\
 \deva{भवतः एकया गर्जनया (हुङ्कारेण) त्रयः लोकाः कम्पन्ते॥}\\

Only you yourself can control your own immense splendor and power.\\
At your single roar, all three worlds tremble.


\vspace{1em}
\newpage
\textbf{\deva{प्रतिपदार्थः} (Meaning of each word):}
\begin{table}[h!]
\centering
\renewcommand{\arraystretch}{1.2}
\begin{tabularx}{\textwidth}{@{} l l X X @{}}
\toprule
\textbf{अवधी} & \textbf{संस्कृतम्} & \textbf{English} & \textbf{Grammar \& Notes} \\
\midrule
\deva{आपन} / \telu{ఆపన} / \textit{āpana}  &  \deva{आत्मीयम्}  &  Own  &   \\
\deva{तेज} / \telu{తేజ} / \textit{teja}  &  \deva{तेजः}  &  Power / Splendor  &   \\
\deva{सम्हारो} / \telu{సమ్హారో} / \textit{samhāro}  &  \deva{सम्भालनम् / नियन्त्रणम्}  &  Control / Handle  &   \\
\deva{आपै} / \telu{ఆపై} / \textit{āpai}  &  \deva{स्वयम् एव}  &  You alone / By yourself  &   \\
\deva{तीनों लोक} / \telu{తీనోం లోక} / \textit{tīnõ loka}  &  \deva{त्रयः लोकाः}  &  Three worlds  &   \\
\deva{हाँक तें} / \telu{హాఁక తేఁ} / \textit{hā̃ka tẽ}  &  \deva{हुङ्कारेण / गर्जनया}  &  By your roar  &   \\
\deva{काँपै} / \telu{కాఁపై} / \textit{kā̃pai}  &  \deva{कम्पन्ते}  &  Tremble / Shiver  &   \\
\bottomrule
\end{tabularx}
\end{table}



\newpage
